\section*{الفصل \ref{ch:b1:alcohol-activation} --- الكحولات: تنشيط مجموعة OH}

\begin{solution}{exo:b1:alcohol-activation:1}
\ce{2CH3CH2OH + 2Na -> 2CH3CH2O- + 2Na+ + H2};
\ce{CH3CH2OH + NaH -> CH3CH2O- + Na+ + H2};
\ce{CH3CH2OH + OH- <=> CH3CH2O- + H2O}، $K = 10^{14.0 - 15.9} = 10^{-1.9}$:
تفاعل جزئي.
\end{solution}

\begin{solution}{exo:b1:alcohol-activation:2}
1-ميثوكسي بروبان؛ ميثوكسي إيثان؛ إيثوكسي بنزين (الفينوكسيد، وهو قاعدة أضعف
من \omterm{def:b1:alcohol-activation:alkoxide}{الألكوكسيد}، يبقى \omterm{def:b1:electronic-effects:nucleophile}{محبًا للنواة} جيدًا).
\end{solution}

\begin{solution}{exo:b1:alcohol-activation:3}
التوزلة في البيريدين تعطي \ce{CH3CH2CH2OTs}؛ ثم
\[ \ce{CH3CH2CH2OTs + I- -> CH3CH2CH2I + TsO-}. \] أما مباشرةً، فكان على اليوديد أن
يطرد \ce{OH-}: فلا تفاعل. وللتوزيلات \omterm{def:b1:electronic-effects:nucleophile}{مجموعة مغادرة} ممتازة،
والظروف ليست حمضية ولا قاعدية بشدة.
\end{solution}

\begin{solution}{exo:b1:alcohol-activation:4}
البوتان-2-أول: بوت-1-ين، وبوت-2-ين-\cip{Z} و-\cip{E}؛ والرئيسي
بوت-2-ين-\cip{E}. 2-ميثيل البنتان-2-أول: 2-ميثيل بنت-2-ين (الرئيسي،
ثلاثي الاستبدال) و2-ميثيل بنت-1-ين.
\end{solution}

\begin{solution}{exo:b1:alcohol-activation:5}
\ce{(CH3)2CHCH2-O-CH2CH3}: 2-ميثيل بروبان-1-أولات الصوديوم مع البروموإيثان،
أو إيثوكسيد الصوديوم مع 1-برومو-2-ميثيل البروبان. والهاليدان كلاهما أوليان،
لكن 1-برومو-2-ميثيل البروبان متفرع بجوار الكربون المتفاعل،
وهذا يبطئ SN2 ويدع \omterm{def:b1:predominant-reaction:strong-acid}{القاعدة القوية} تُحدث حذفًا (2-ميثيل البروبين).
والاختيار الأول أفضل.
\end{solution}

\begin{solution}{exo:b1:alcohol-activation:6}
تحفظ \omterm{def:b1:alcohol-activation:sulfonate}{التوزيلات} الترتيب \cip{R}؛ ويقلبه SN2 مع الإيثوكسيد:
2-إيثوكسي أوكتان-\cip{S} (تحتل OEt الرتبة الأولى كما كانت OTs). ومع حمض الكبريتيك
والحرارة: حذف إلى أوكتينات (أساسًا أوكت-2-ين-\cip{E})، عبر \omterm{def:b1:electronic-effects:intermediates}{كاتيون
كربوني} ثانوي، مع ضياع الكيمياء الفراغية.
\end{solution}

\begin{solution}{exo:b1:alcohol-activation:7}
برتنة OH؛ ثم فقد الماء نحو \omterm{def:b1:electronic-effects:intermediates}{الكاتيون الكربوني} الثالثي (بطيء)؛ ثم
أسره بالأيون \ce{Cl-}: 2-كلورو-2-ميثيل البروبان. وهو سريع لأن \omterm{def:b1:electronic-effects:intermediates}{الكاتيون الكربوني}
الثالثي ثابت، والكلوريد، غير الذواب في الحمض، ينفصل.
\end{solution}

\begin{solution}{exo:b1:alcohol-activation:8}
\ce{CH3CH2OH + H+ -> CH3CH2OH2+}؛ ويهاجم جزيء إيثانول ثانٍ الكربون
من الخلف، بينما يغادر الماء (SN2): \ce{(CH3CH2)2OH+}، الذي يفقد بروتونًا:
الإيثوكسي إيثان. والتفاعل المنافس: تأخذ قاعدة (الإيثانول، \ce{HSO4-}) بروتون $\beta$
من \ce{CH3CH2OH2+} بينما يغادر الماء (E2): الإيثين.
\end{solution}

\begin{solution}{exo:b1:alcohol-activation:9}
2-ميثيل البروبان-2-أول (ثالثي)، فيعطي 2-كلورو-2-ميثيل البروبان، غير الذواب.
\end{solution}

\begin{solution}{exo:b1:alcohol-activation:10}
تعطي البرتنة وفقد الماء \omterm{def:b1:electronic-effects:intermediates}{كاتيونًا كربونيًا} ثانويًا بجوار
كربون رباعي. فتنتقل مجموعة ميثيل من ذلك الكربون مع زوجها الرابط
إلى الكربون الكاتيوني: فتصير الشحنة الموجبة على \omterm{def:b1:electronic-effects:carbon-class}{كربون
ثالثي}، أكثر ثباتًا. ويعطي فقد بروتون من كربون مجاور
2,3-ثنائي ميثيل بوت-2-ين الرباعي الاستبدال، وهو الألكين الأكثر ثباتًا.
\end{solution}

\begin{solution}{exo:b1:alcohol-activation:11}
من البوتان-2-أول-\cip{R}: \omterm{def:b1:alcohol-activation:sulfonate}{التوزيلات} (\cip{R})، ثم ميثوكسيد الصوديوم (SN2،
انقلاب): 2-ميثوكسي بوتان-\cip{S}. ومن البوتان-2-أول-\cip{S}: حضّر
\omterm{def:b1:alcohol-activation:alkoxide}{الألكوكسيد} (NaH) وأضف اليودوميثان: فالرابطة C--O \omterm{def:b1:stereochemistry-in-depth:stereocentre}{للمركز الفراغي} لا
تُمس، فنحصل على 2-ميثوكسي بوتان-\cip{S} مرة أخرى.
\end{solution}

\begin{solution}{exo:b1:alcohol-activation:12}
$Q = [\text{إيثر}][\ce{H2O}]/[\ce{EtOH}]^2$. وإزالة الماء كلما تكوّن
تُبقي $[\ce{H2O}]$ و$Q$ صغيرين: $Q < K$ فيستمر التفاعل في تكوين
الإيثر حتى يُستهلك الكحول.
\end{solution}

\begin{solution}{pb:b1:alcohol-activation:1}
\textbf{1.} ثالثي بوتوكسيد الصوديوم مع هاليد ميثيل؛ وميثوكسيد الصوديوم مع
هاليد ثالثي البوتيل.
\textbf{2.} الثاني: فالميثوكسيد، وهو \omterm{def:b1:predominant-reaction:strong-acid}{قاعدة قوية}، يُحدث حذفًا من
الهاليد الثالثي، \ce{(CH3)3CBr + CH3O- -> (CH3)2C=CH2 + CH3OH + Br-}.
\textbf{3.} \ce{(CH3)3CO- + CH3I -> (CH3)3COCH3 + I-}: يهاجم أكسجين \omterm{def:b1:alcohol-activation:alkoxide}{الألكوكسيد}
كربون الميثيل من الخلف بينما يغادر اليوديد (SN2).
\textbf{4.} بالصوديوم (أو بهيدريد الصوديوم): \ce{2(CH3)3COH + 2Na ->
2(CH3)3CO- + 2Na+ + H2}. والهيدروكسيد قاعدة أضعف من اللازم: $K = 10^{14.0 -
19.2} = 10^{-5.2}$.
\textbf{5.} كان الماء سيبرتن \omterm{def:b1:alcohol-activation:alkoxide}{الألكوكسيد} ويهاجم الهاليد.
\textbf{6.} لا: لا \omterm{def:b1:stereochemistry-in-depth:stereocentre}{مركز فراغي} فيه.
\textbf{7.} \babelsublr{\foreignlanguage{arabic}{بوتان-2-أول-\cip{R} + TsCl (بيريدين)} $\to$ \foreignlanguage{arabic}{\omterm{def:b1:alcohol-activation:sulfonate}{توزيلات} البوتان-2-يل-\cip{R}}}:
احتفاظ.
\textbf{8.} 2-ميثوكسي بوتان-\cip{S}، بانقلاب SN2.
\textbf{9.} بوتان-2-أول-\cip{S}.
\textbf{10.} 2-ميثوكسي بوتان-\cip{R}: تُحفظ الرابطة C--O للكربون
الفراغي المركز؛ ولا تتكوّن إلا الرابطة O--C(ميثيل).
\textbf{11.} E2 (فالميثوكسيد \omterm{def:b1:predominant-reaction:strong-acid}{قاعدة قوية} \omterm{def:b1:electronic-effects:carbon-class}{والكربون ثانوي})،
فتتكوّن بوتينات؛ وتحدّ منه القاعدة الصغيرة، \omterm{def:b1:electronic-effects:nucleophile}{والمجموعة المغادرة} الممتازة
ودرجة الحرارة المعتدلة.
\textbf{12.} يتأين الكحول المتبرتن إلى \omterm{def:b1:electronic-effects:intermediates}{كاتيون كربوني} ثانوي مستوٍ،
يأسره الميثانول على الوجهين: إيثر راسيمي، مع
بوتينات.
\textbf{13.} برتنة OH، ثم فقد الماء (بطيء) نحو \omterm{def:b1:electronic-effects:intermediates}{الكاتيون الكربوني}
الثالثي، ثم فقد بروتون $\beta$.
\textbf{14.} 2-ميثيل بوت-2-ين (الرئيسي، ثلاثي الاستبدال)
و2-ميثيل بوت-1-ين.
\textbf{15.} \omterm{def:b1:electronic-effects:intermediates}{الكاتيون الكربوني} الثالثي معاق ويفقد بروتونًا بسهولة
أكبر بكثير مما يستطيع جزيء كحول آخر أن يهاجمه.
\textbf{16.} لإزالة ناتج: فمع إبقاء $Q$ صغيرًا يستمر التفاعل،
ولا يبقى الألكين على تماس مع الحمض.
\textbf{17.} حاجز أول مرتفع (فقد الماء، المحدد للسرعة) نحو
\omterm{def:b1:electronic-effects:intermediates}{الكاتيون الكربوني}، ثم حاجز صغير نحو الألكين.
\textbf{18.} لا: \omterm{def:b1:electronic-effects:intermediates}{الكاتيون الكربوني} الأولي غير ثابت إلى حد كبير؛ والكحولات الأولية
يُنزع منها الماء \omterm{def:b1:elimination:elimination}{بالآلية E2} على الكحول المتبرتن، عند درجة حرارة أعلى.
\textbf{19.} $10.0/74.0 = \qty{0.135}{mol}$.
\textbf{20.} $1.10 \times 0.135 \times 141.9 = \qty{21.1}{g}$.
\textbf{21.} $0.135 \times 88.0 = \qty{11.9}{g}$.
\textbf{22.} $0.75 \times 11.9 =$ \textbf{\qty{8.9}{g} من MTBE}.
\end{solution}
